% !TEX program = xelatex
% ============================================================
%  EduChain：基于区块链的智慧教育学术诚信保障系统研究报告
%  仿《计算机学报》(Chinese Journal of Computers) 排版风格
%  编译方式：xelatex polaris.tex （需两次以生成目录与交叉引用）
%  依赖：ctex 宏集（TeX Live / MiKTeX 自带）
% ============================================================
\documentclass[10pt,twocolumn]{ctexart}

% ---------- 页面与版式 ----------
\usepackage[a4paper,top=2.0cm,bottom=2.0cm,left=1.8cm,right=1.8cm,columnsep=0.7cm]{geometry}
\usepackage{indentfirst}
\usepackage{multicol}
\usepackage{titlesec}
\usepackage{fancyhdr}

% ---------- 数学与符号 ----------
\usepackage{amsmath,amssymb,amsthm}
\usepackage{bm}

% ---------- 表格 ----------
\usepackage{booktabs}
\usepackage{array}
\usepackage{tabularx}
\usepackage{multirow}

% ---------- 图形 ----------
\usepackage{graphicx}

% ---------- 算法与代码 ----------
\usepackage{listings}
\usepackage{xcolor}
\usepackage[ruled,linesnumbered]{algorithm2e}

% ---------- 超链接 ----------
\usepackage[hidelinks]{hyperref}

% ---------- 代码样式 ----------
\lstset{
  basicstyle=\ttfamily\small,
  keywordstyle=\color{blue!70!black}\bfseries,
  commentstyle=\color{green!45!black}\itshape,
  stringstyle=\color{red!60!black},
  numbers=left,
  numberstyle=\tiny\color{gray},
  numbersep=6pt,
  frame=single,
  framesep=4pt,
  breaklines=true,
  showstringspaces=false,
  language=Python,
  tabsize=2,
  columns=fullflexible
}

% ---------- 定理类环境 ----------
\newtheorem{theorem}{定理}
\newtheorem{lemma}{引理}
\newtheorem{definition}{定义}

% ---------- 章节标题样式 ----------
\titleformat{\section}{\zihao{4}\heiti}{\thesection}{0.8em}{}
\titleformat{\subsection}{\zihao{5}\heiti}{\thesubsection}{0.6em}{}
\titleformat{\subsubsection}{\zihao{5}\songti\bfseries}{\thesubsubsection}{0.6em}{}

% ---------- 页眉页脚 ----------
\pagestyle{fancy}
\fancyhf{}
\fancyhead[C]{\zihao{-5}\songti EduChain：基于区块链的智慧教育学术诚信保障系统}
\fancyfoot[C]{\thepage}
\renewcommand{\headrulewidth}{0.4pt}

\begin{document}

% ============================================================
%  标题区（跨栏）
% ============================================================
\twocolumn[
  \begin{@twocolumnfalse}
  \begin{center}
    {\zihao{2}\heiti EduChain：基于区块链的智慧教育\\学术诚信保障系统研究报告\par}
    \vspace{1.2em}
    {\zihao{4}\kaishu 研究报告\par}
    \vspace{1.5em}
  \end{center}
  \end{@twocolumnfalse}
]

% ============================================================
%  摘要与关键词
% ============================================================
\noindent{\heiti\zihao{-4} 摘\quad 要}\quad
针对智慧教育场景中学生作业评价存在的主观性强、抄袭与代写难以追溯、评价记录易被篡改等结构性问题，本文设计并实现了基于区块链的智慧教育学术诚信保障系统 EduChain。系统将"乐观默认、异常裁决"的协议结构引入学术评价流程：默认信任学生提交的作业是原创且合格的、AI 与同行评审的结论是可靠的，仅在出现学术争议时才触发代价较高的链上仲裁。系统采用链上协议层与链下执行层分离的双层架构，链上承载内容承诺、状态转移与最终裁决，链下由 AI 评审代理与隔离沙箱提供语义分析与可执行验证。通过学分质押-学术声誉双资产博弈论模型，从理论上证明了诚实行为构成纳什均衡。原型系统的端到端验证表明，EduChain 具备工程可行性，能够为智慧教育提供可信的学术诚信数据管理基础设施。

\vspace{0.6em}
\noindent{\heiti\zihao{-4} 关键词}\quad
智慧教育；区块链数据管理；乐观执行；博弈论激励；学术诚信；去中心化仲裁

\vspace{1.0em}

% ============================================================
%  正文
% ============================================================
\section{研究背景与问题}

在智慧教育与计算教育学场景中，学生作业（尤其是编程作业）的评价长期面临三类结构性难题。其一，评分主观性强，不同评审者对同一作业的评分可能差异显著，缺乏客观、统一的衡量标准。其二，抄袭与代写难以追溯，传统的查重系统依赖中心化数据库，既容易被绕过或篡改，也难以对检测结果进行独立验证。其三，评价记录易被篡改或丢失，集中式教务系统将所有信任集中于单一机构，既缺乏过程透明性，也难以支撑跨校、跨学期的学术档案互认。

区块链技术的不可篡改、去中心化与可追溯特性为解决上述问题提供了新的技术路径。近年来，教育区块链领域的研究主要聚焦于学籍档案管理与学历证书存证层面，但将区块链技术系统性地应用于日常作业评价与学术争议解决的工作尚不多见。

在区块链协议设计层面，Optimistic Rollup 的乐观执行思想提供了重要启发：默认假设交易有效，仅在争议时触发验证。在去中心化争议解决方面，Kleros 的去中心化仲裁协议采用 Schelling Point 博弈激励陪审员投票，但其依赖纯主观投票、缺乏客观的技术验证手段。基于大语言模型（LLM）的代码评审方法展现了理解程序语义的潜力，但幻觉问题导致评审结论的可靠性存疑。

本文提出的 EduChain 系统旨在将"乐观默认、异常裁决"的协议结构引入学术评价流程，结合密码学承诺、博弈论激励与可重放验证沙箱，构建一套去中心化、可追责、可经济惩罚、且任何第三方可独立复现的学术诚信保障协议。

\section{系统设计}

\subsection{系统模型与参与方}

EduChain 系统中存在五类参与方：教师（任务发布方）、学生（作业提交者）、同行评议者（独立评分人）、举报者（争议提出者）与仲裁者（委员会成员）。除教师外，其余四类均为注册参与者，仅通过链上协议交互。

威胁假设上，系统不依赖任何中心化可信第三方，假设任意参与方均可能为理性自利主体，包括抄袭代写、敷衍评审、串通打分、恶意诬告等。在此模型下，协议的核心目标是使诚实行为（原创提交、认真评审、据实举报）成为所有参与方的占优策略。

\subsection{双层架构设计}

EduChain 采用链上协议层与链下执行层分离的双层架构。

\textbf{链上协议层}由四个 Solidity 智能合约组成：Registry 合约管理去中心化身份注册、学分质押锁定与学术声誉追踪，是抗女巫攻击的第一道防线；DisputeResolution 合约承载完整的五阶段作业评审生命周期，统一调度任务发布、作业提交、同行评审、争议发起、乐观结算与仲裁结算；ArbitrationCommittee 合约实现委员会的可验证随机选举与 EIP-712 签名聚合验证；StakeOracle 合约根据可调的教育激励参数动态计算最低质押门槛。

\textbf{链下执行层}包含三个核心服务：AI 评审代理承担作业质量评审、抄袭检测、验证代码生成与仲裁推理等需要语义理解的任务；隔离沙箱服务为查重与功能验证代码提供完全隔离、确定性可重现的执行环境；去中心化存储网络（IPFS）承载完整的作业包、验证证据与评审记录。

\subsection{五阶段学术评审流水线}

流水线状态机由六种状态组成：开放（Open）、已提交（Proposed）、评审中（InReview）、被争议（Challenged）、已最终化（Finalized）、已罚没（Slashed）。

\textbf{阶段一：作业发布。} 教师发布作业任务，附带要求哈希、评分标准（硬约束）、争议窗口时长（默认48小时）、最低质押与作为奖励池的学分。

\textbf{阶段二：学生提交与质押。} 学生将作业的三棵 Merkle 树根（状态根承载内容哈希、证据根承载测试结果等证据、追踪根承载开发过程轨迹）连同完整作业包的 IPFS 指针提交至链上，同时锁定一笔学分质押作为对原创性的经济担保。

\textbf{阶段三：AI 评审与同行互评。} AI 评审代理首先对作业进行多维自动评审。同时多位同行评议者采用提交-揭示两阶段机制独立打分：提交阶段上链分数哈希，揭示阶段公开分数与盐值，防止"跟随串谋"。

\textbf{阶段四：争议窗口。} 任何注册参与者均可在窗口期内发起争议。争议的核心载荷是一份可执行的验证代码——例如证明代码与公开来源高度相似的查重对比，或证明功能不符合声明的功能测试，并附以举报者质押。

\textbf{阶段五甲：乐观结算路径。} 若窗口期无人争议且同行评分均值达到通过阈值，协议进入乐观结算：学生取回质押并获得学分奖励、声誉+10；同行评议者按 Shapley 值近似分配评审奖励。

\textbf{阶段五乙：仲裁结算路径。} 若发起争议，协议进入仲裁路径：VRF 从声誉 $\geq$ 200 的活跃参与者中随机选出仲裁委员会 → 委员在隔离沙箱中独立运行举报者的验证代码 → 对判定结果进行 EIP-712 签名投票（67\% 法定人数）→ 链上原子执行学分罚没与声誉调整。若抄袭成立，学生质押被罚没并按 $\alpha$ 比例分配给举报者，学生声誉-50、举报者+20。若争议驳回，举报者质押被对称罚没，学生获补偿，学生声誉+10、举报者-30。

\subsection{博弈论激励机制}

将参与者交互建模为不完全信息博弈。学生策略集 $\Sigma_s = \{\text{诚实提交}, \text{作弊}\}$，举报者策略集 $\Sigma_c = \{\text{据实举报}, \text{恶意诬告}, \text{不举报}\}$。

设 $C_h$ 为诚实完成作业成本，$C_c$ 为作弊成本，$S_p$ 为学生质押额，$R$ 为学分奖励，$P_{detect}$ 为抄袭被检测概率，$P_{arb}$ 为仲裁正确率，$\alpha$ 为罚没学分分配给举报者的比例。

诚实学生期望收益 $E[\pi_s^{honest}] = R - C_h$。作弊学生期望收益 $E[\pi_s^{cheat}] = (R - C_c)(1 - P_{detect} \cdot P_{arb}) - P_{detect} \cdot P_{arb} \cdot S_p$。令诚实占优，解得最低质押约束：
\begin{equation}
S_p > \frac{C_h - C_c}{P_{detect} \times P_{arb}}
\end{equation}

基于上式设计链上动态质押预言机，根据系统实时参数自动计算最低质押额。

\begin{theorem}
当质押参数满足式(1)时，（诚实提交，据实举报）构成纳什均衡。
\end{theorem}

代入典型参数（$P_{detect}=0.7$, $P_{arb}=0.95$, $C_h=2.0$学分, $C_c=0.1$学分, $S_p=2.86$学分），作弊学生期望收益为 $-0.065$ 学分 $< 0$，恶意举报者期望收益为 $-0.864$ 学分 $< 0$，据实举报者期望收益为 $1.0$ 学分 $> 0$。诚实行为构成占优策略。

此外，系统采用质押-声誉双资产模型防止女巫攻击：注册需学分质押 + W3C DID 身份认证，声誉初始100，仲裁候选需声誉 $\geq$ 200，参与者综合权重 $=$ 质押 $\times$ 声誉，声誉无法靠资金快速购买，必须通过持续诚实行为缓慢积累。

\subsection{AI 评审代理与沙箱验证}

AI 评审代理由大语言模型驱动，依次完成作业用途识别、正确性检查、代码质量评估、抄袭指标检测与评分标准对照。审计深度依据激励动态调整：作业分值高时增加分析轮次与搜索遍数，分值低时降低强度，但对硬性要求的核验始终保留为质量下限。

举报者代理在定位问题后，通过"生成—编译—修复"闭环将问题翻译为可执行验证测试：生成初版代码 → 编译捕获错误 → 回传 LLM 修复 → 循环直至通过。该闭环使产出从"可读"升级为"可验证"。

隔离沙箱遵循网络隔离、资源约束、文件系统受限三类硬约束。判定机械化：编译失败裁定争议不成立；测试通过且断言成立（如相似度超阈值）裁定争议成立。执行输出计算哈希得到不可伪造的轨迹哈希。

\section{系统实现}

\subsection{原型系统架构}

EduChain 原型系统采用分层解耦架构，自底向上划分为链上协议层、链下服务层与用户交互层。

链上协议层基于 Solidity 0.8.24 实现，合约通过 Foundry 框架编译、测试与部署。链下服务层基于 Python 3.10+ / FastAPI 实现，封装 AI 评审、代码查重验证、沙箱执行与仲裁辅助评估服务。用户交互层基于 React + TypeScript + Tailwind CSS 实现响应式 Web 界面，通过 ethers.js 与链上合约交互。

\subsection{核心 API 接口}

系统提供以下 REST API 端点：\texttt{POST /api/audit/perform} 执行 AI 自动评审，输入学生代码与评分约束，返回评审状态根、证据根与质量评分；\texttt{POST /api/poc/generate} 生成抄袭验证代码，输入问题类型与目标代码，返回可编译的验证测试；\texttt{POST /api/sandbox/replay} 执行沙箱验证，输入验证代码与目标源码，返回判定结论与重放轨迹哈希；\texttt{POST /api/arbitration/evaluate} 执行仲裁评估，输入争议双方材料与沙箱结果，返回结构化裁决与置信度。

\section{实验评估}

\subsection{实验设置}

实验在配备 LLM API（DeepSeek）的环境中进行，数据集包含 16 份学生代码提交（8份原创 + 8份存在问题，包括抄袭、逻辑错误、代码风格问题）。基线方法包括：B1（单 Agent LLM 评审）、B2（多 Agent 投票）、B3（传统静态分析工具）、B4（全量验证，无乐观机制）。

\subsection{核心实验结果}

实验结果表明，EduChain 在关键指标上表现优异。在安全性能方面，EduChain 在保持最高召回率（0.900）的同时，通过沙箱验证将误报率降至 0.125，F1 分数达 0.865。在效率方面，乐观路径仅消耗 450K gas/task，相较于全量验证（1,180K gas/task）节省约 62\%。在博弈论验证方面，诚实学生期望收益为正值（1.800 ETH），作弊学生与恶意举报者期望收益均为负值，验证了激励相容设计的有效性。

\section{结论}

针对智慧教育场景中学生作业评价的三大结构性难题，本文提出并实现了基于区块链的 EduChain 学术诚信保障系统。核心贡献包括：将乐观执行思想创新性地应用于教育评价场景，实现按需验证的高效机制；构建学分质押-学术声誉双资产博弈论模型，从理论上证明诚实行为构成纳什均衡；设计基于 VRF 与结构化签名投票的去中心化仲裁协议，以可执行验证替代主观投票；实现 LLM 驱动的 AI 评审代理与确定性沙箱重放机制。原型系统的端到端验证表明，EduChain 具备工程可行性，为智慧教育提供了可信的去中心化学术诚信数据管理基础设施。

未来工作将在大规模真实教学场景中进行定量实验验证，扩展系统对多模态作业类型的支持，以及探索基于零知识证明的隐私保护评审方案。

\end{document}
